\documentclass{article}

\usepackage{amsmath, amsthm, amssymb, amsfonts}
\usepackage{thmtools}
\usepackage{graphicx}
\usepackage{setspace}
\usepackage{geometry}
\usepackage{float}
\usepackage[backend=biber,style=ieee,sorting=ynt]{biblatex}
\usepackage[hidelinks]{hyperref}
\usepackage[utf8]{inputenc}
\usepackage[spanish]{babel}
\usepackage{framed}
\usepackage[dvipsnames]{xcolor}
\usepackage{tcolorbox}
\usepackage{tikz}
\usetikzlibrary{positioning}
\usetikzlibrary{arrows.meta}
\usepackage{caption}
\usepackage{longtable}
\usepackage{pdflscape}
\usepackage{svg}
\usepackage{subcaption}
\usepackage{caption}
\usepackage{multirow}
\usepackage{listings}
\usepackage{cancel}
\usepackage{fancyhdr}
\usepackage{booktabs}
\usepackage{array}
\usepackage{ragged2e}

\addbibresource{referencias.bib}


\colorlet{LightGray}{White!90!Periwinkle}
\colorlet{LightOrange}{Orange!15}
\colorlet{LightGreen}{Green!15}



\newcommand{\HRule}[1]{\rule{\linewidth}{#1}}

\declaretheoremstyle[name=Theorem,]{thmsty}
\declaretheorem[style=thmsty,numberwithin=section]{theorem}
\tcolorboxenvironment{theorem}{colback=LightGray}

\declaretheoremstyle[name=Proposition,]{prosty}
\declaretheorem[style=prosty,numberlike=theorem]{proposition}
\tcolorboxenvironment{proposition}{colback=LightOrange}

\declaretheoremstyle[name=Principle,]{prcpsty}
\declaretheorem[style=prcpsty,numberlike=theorem]{principle}
\tcolorboxenvironment{principle}{colback=LightGreen}

\newcolumntype{L}[1]{>{\raggedleft\let\newline\\\arraybackslash\hspace{0pt}}m{#1}}
\newcolumntype{C}[1]{>{\centering\let\newline\\\arraybackslash\hspace{0pt}}m{#1}}
\newcolumntype{R}[1]{>{\raggedright\let\newline\\\arraybackslash\hspace{0pt}}m{#1}}

\hypersetup{
    pdftitle={Avance De Proyecto II: Formulación Del Problema Como Agente Inteligente},
    pdfauthor={Alvarado Ludwig, Pinilla Kevin, Mahecha Alberto, Ospina María},
    pdfsubject={Avance proyecto},
    pdfkeywords={proyecto, agentes},
    pdfcreator={GNU Emacs desde Arch GNU/Linux (btw) y LaTeX con hyperref},
    pdfproducer={pdflatex}
}


\setstretch{1.2}
\geometry{
    textheight=9in,
    textwidth=5.5in,
    top=1in,
    headheight=12pt,
    headsep=25pt,
    footskip=30pt
}

\lstdefinestyle{bashstyle}{
    language=bash,
    basicstyle=\ttfamily,
    backgroundcolor=\color{gray!10},
    keywordstyle=\color{blue},
    commentstyle=\color{green!40!black},
    stringstyle=\color{red},
    showstringspaces=false,
    numbers=left,
    numberstyle=\tiny\color{gray},
    breaklines=true,
    breakatwhitespace=true,
    frame=tb,
    rulecolor=\color{black!70},
    framerule=0.5pt,
    tabsize=4,
    captionpos=b
}

\lstdefinestyle{javastyle}{
    language=Java,
    basicstyle=\ttfamily,
    backgroundcolor=\color{gray!10},
    keywordstyle=\color{blue},
    commentstyle=\color{green!40!black},
    stringstyle=\color{red},
    showstringspaces=false,
    numbers=left,
    numberstyle=\tiny\color{gray},
    breaklines=true,
    breakatwhitespace=true,
    frame=tb,
    rulecolor=\color{black!70},
    framerule=0.5pt,
    tabsize=4,
    captionpos=b
}

\lstdefinestyle{pythonstyle}{
    language=Python,
    basicstyle=\ttfamily,
    backgroundcolor=\color{gray!10},
    keywordstyle=\color{blue},
    commentstyle=\color{green!40!black},
    stringstyle=\color{red},
    showstringspaces=false,
    numbers=left,
    numberstyle=\tiny\color{gray},
    breaklines=true,
    breakatwhitespace=true,
    frame=tb,
    rulecolor=\color{black!70},
    framerule=0.5pt,
    tabsize=4,
    captionpos=b
}

\lstdefinestyle{cppstyle}{
    language=C++,
    basicstyle=\ttfamily,
    backgroundcolor=\color{gray!10},
    keywordstyle=\color{blue},
    commentstyle=\color{green!40!black},
    stringstyle=\color{red},
    showstringspaces=false,
    numbers=left,
    numberstyle=\tiny\color{gray},
    breaklines=true,
    breakatwhitespace=true,
    frame=tb,
    rulecolor=\color{black!70},
    framerule=0.5pt,
    tabsize=4,
    captionpos=b
}



\pagestyle{fancy}
\fancyhf{}
\renewcommand{\headrulewidth}{0pt}
\fancyfoot[C]{\thepage}
\fancyfoot[C]{\footnotesize \thepage \quad \textbar{} \quad Este trabajo está bajo una licencia CC BY-SA 4.0.\\ Más info: \url{https://creativecommons.org/licenses/by/4.0/}}


% ------------------------------------------------------------------------------

\begin{document}

% ------------------------------------------------------------------------------
% Cover Page and ToC
% ------------------------------------------------------------------------------

\title{ \normalsize \textsc{}
	\\ [2.0cm]
	\HRule{1.5pt} \\
	\LARGE \textbf{\uppercase{FORMULACIÓN DEL PROBLEMA COMO AGENTE INTELIGENTE}
		\HRule{2.0pt} \\ [0.6cm] \LARGE{Universidad de Bogotá Jorge Tadeo Lozano} \vspace*{11\baselineskip}}
}
\date{}
\author{\textbf{Alvarado Becerra Ludwig} \\
  \textbf{Pinilla Naranjo Kevin} \\
  \textbf{Mahecha Reyes Alberto} \\
  \textbf{Ospina Rocha María} \\ 
	 Inteligencia Artificial - 2026-2S}

\maketitle
\thispagestyle{empty}
\newpage

\tableofcontents
\listoftables
\thispagestyle{empty}
\newpage
\setcounter{page}{1}

\section{Descripción y Justificación del Problema}

El sistema de transporte masivo TransMilenio en Bogotá enfrenta severos problemas de congestión durante las horas pico, caracterizados por andenes colapsados en estaciones clave y tiempos de espera impredecibles para los usuarios~\cite{newell1982queueing,transmilenio2026datos}. Según cifras oficiales de TransMilenio, el sistema moviliza más de 2 millones de viajes diarios~\cite{transmilenio2026datos}.

La causa principal de esta situación radica en la rigidez de la programación de despachos, la cual genera una asimetría de ocupación: mientras las rutas corrientes experimentan saturación severa dejando pasajeros varados en plataforma, los servicios expresos que transitan de forma paralela operan con capacidad ociosa~\cite{mobilitydata2026gtfs,transmilenio2026datos}.

La adquisición de nuevos vehículos no es viable a corto plazo debido a su elevado costo de capital (aproximadamente \$1.5 millones USD por bus) y a la saturación de la capacidad física en los patios operativos~\cite{russell2020artificial}. Por lo tanto, la alternativa viable consiste en formular un Agente Inteligente capaz de optimizar y redistribuir los intervalos de despacho sobre la flota existente, haciendo coincidir la oferta de transporte con los requerimientos reales de la demanda~\cite{russell2020artificial}.

\section{Definición Del Agente Inteligente}

Se define como agente al \textbf{Despachador Racional de Frecuencias Troncales (FlowTM-Scheduler)}.

\begin{itemize}
  \item \textbf{Entidad Operativa:} No se confunde con la interfaz gráfica ni con el motor de almacenamiento. El agente es el software decisor que evalúa el estado del corredor y determina la asignación algorítmica de frecuencias de paso para las rutas en operación.
  \item \textbf{Decisión que toma:} El agente decide el Vector de Intervalos de Despacho \(H = \{h_r\}\) para cada ruta \(r \in {B13, B18, B23, B27, B74, B75}\) en minutos, restringido por la flota finita disponible:

        Decisión: \(H = \{ h_{B13}, h_{B18}, h_{B23}, h_{B27}, h_{B74}, h_{B75}\} \in [2.0, 15.0]^6 \subset \mathbb{R}^6\)

        Sujeto a: \(\sum_r B_{\mathrm{usados}}(h_r) \leq B_{\mathrm{max}}\). Donde: \(B_{\mathrm{max}}\) es una constante, la flota física acotada, en total 106 buses, \(B_{usados}\) se define como:

        \[
        B_{\mathrm{usados}}(h_r) = \dfrac{t_{\mathrm{sim}}}{h_r}
        \]

        \(t_{\mathrm{sim}}\) es el tiempo en minutos que dura la simulación.

\end{itemize}

\section{Modelo PEAS/REAS}

\begin{longtable}{|>{\bfseries}p{0.18\textwidth} |p{0.76\textwidth}|}
    \caption{Especificación PEAS del agente inteligente FlowTM}
    \label{tab:peas-flowtm} \\

    \toprule
    \textbf{Dimensión PEAS} & \textbf{Descripción específica del proyecto FlowTM} \\
    \midrule
    \endfirsthead

    \multicolumn{2}{c}%
    {{\tablename\ \thetable{} -- continuación}} \\
    \toprule
    \textbf{Dimensión PEAS} & \textbf{Descripción específica del proyecto FlowTM} \\
    \midrule
    \endhead

    \midrule
    \multicolumn{2}{r}{{Continúa en la siguiente página}} \\
    \endfoot

    \bottomrule
    \endlastfoot

    \textbf{Rendimiento (Performance)}
    &
    Función de utilidad multicriterio \(U(S)\) que penaliza la ineficiencia
    del sistema:

    \[
        U(S) =
        -\left[
        w_1 \bar{W}
        + w_2(\mathrm{Sobrecupo}_{95\%})^{1.5}
        + w_3 Q_{\mathrm{remanente}}
        + w_4 B_{\mathrm{usados}}
        \right]
    \]

    donde \(\bar{W}\) es el tiempo de espera medio ponderado (min),
    \(\mathrm{Sobrecupo}_{95\%}\) es la fracción de buses con ocupación
    \(\geq 95\%\), penalizada exponencialmente con exponente \(1.5\) para
    castigar el hacinamiento extremo, \(Q_{\mathrm{remanente}}\) son los
    pasajeros que quedan en plataforma al terminar la franja y
    \(B_{\mathrm{usados}}\) es la flota consumida. Los pesos son
    \(w_1=0.40\), \(w_2=0.30\), \(w_3=0.15\) y \(w_4=0.15\).
    \\ \hline 

    Entorno (Environment)
    &
    Corredor troncal Autonorte--(Calle 45 Caracas), en sentido sur a
    norte, compuesto por 29 estaciones secuenciales, carriles
    exclusivos de tráfico segregado, 6 servicios troncales
    (B13, B18, B23, B27, B74, B75), andenes (puertas de embarque) con
    capacidad de parada acotada (máximo 2 bahías físicas de
    estacionamiento simultáneas por estación) y una afluencia acumulada
    de 45\,000 pasajeros en la franja pico.
    \\ \hline

    Actuadores (Actuators)
    &
    Módulo exportador de especificación oficial de tránsito GTFS
    (\texttt{frequencies.txt}), el cual publica los campos estructurados
    (\texttt{route\_id}, \texttt{start\_time}, \texttt{end\_time},
    \texttt{headway\_secs}, \texttt{exact\_times=0}) para ser consumidos
    por los sistemas de gestión y despacho de flota (SAE).
    \\ \hline

    Sensores (Sensors)
    &
    \begin{enumerate}
        \item Ingesta de archivos \texttt{.Parquet} de la matriz de
        validaciones de tarjetas tullave
        (\(\lambda(s,t)\) agregada en intervalos de 15 min).

        \item Conteos de desocupación en torniquetes externos
        (\textit{Salidas\_S}) para el cálculo de probabilidades de
        descenso Origen--Destino.

        \item Topología georreferenciada del GeoJSON oficial
        (latitudes/longitudes de paradas) y archivos GTFS
        (\texttt{stops.txt}, \texttt{trips.txt},
        \texttt{stop\_times.txt}).
    \end{enumerate}
    \\

\end{longtable}

\section{Caracterización del Entorno}

\begin{itemize}
  \item \textbf{Parcialmente Observable:} El agente percibe las validaciones de acceso al andén y las salidas por torniquete, pero desconoce el destino final de cada usuario individual al momento de entrar y no observa la velocidad instantánea de los buses entre estaciones afectadas por semáforos [4].

  \item \textbf{Estocástico:} Las llegadas de pasajeros se rigen por un Proceso de Poisson No Homogéneo (NHPP) y los tiempos de abordaje (Dwell Time) varían aleatoriamente en función de la aglomeración de las puertas.

  \item \textbf{Dinámico:} La demanda y la ocupación de los buses evolucionan continuamente mientras el agente evalúa y calcula alternativas de despacho.

  \item \textbf{Secuencial:} Las frecuencias asignadas a las 06:00 a.m. condicionan directamente el tamaño de las colas y la disponibilidad de flota a las 07:30 a.m.

  \item \textbf{Continuo:} El tiempo físico y el desplazamiento espacial de los buses fluyen continuamente, requiriendo su integración numérica en SimPy.

  \item \textbf{Individual / Centralizado:} Existe un único agente decisor global responsable de optimizar el corredor completo.

\end{itemize}


\section{Ciclo Percepción -– Decisión –- Acción}

El proceso continuo del agente se estructura en 5 pasos fundamentales:

\begin{enumerate}
  \item Paso 1 (Percibir): El agente lee las entradas de usuarios por torniquete en cada estación, las salidas registradas y el número de buses disponibles en patio.
  \item Paso 2 (Simular): Prueba internamente un esquema de horarios propuesto y proyecta la acumulación de filas en andén, la ocupación de los buses y el tiempo de parada por estación mediante balance de flujos.
  \item Paso 3 (Evaluar): Calcula la puntuación del esquema simulado comparándolo contra el horario base mediante la función de utilidad U(S).
  \item Paso 4 (Decidir): Selecciona la configuración de intervalos que ofrece el mejor balance entre menor tiempo de espera y menor sobrecupo, sin sobrepasar el límite de 106 buses.
  \item Paso 5 (Actuar): Genera y exporta la tabla final de programación en formato GTFS (frequencies.txt) para su ejecución operativa en la flota.
\end{enumerate}


\section{Datos Requeridos y Viabilidad Arquitectónica}

\subsection{Inventario y Acceso Real a Fuentes de Datos}




\addcontentsline{toc}{section}{Referencias}
\printbibliography

\newpage

\thispagestyle{empty}

\vspace*{0.3\textwidth}

\begin{figure}[H]
  \centering
  \includesvg[width=\textwidth]{img/LogoUtadeo.svg}
\end{figure}



\end{document}
